% GENERATED FILE. Edit canonical lessons, then run npm run generate:books.
% canonical-source-hash: fnv1a64:c21f06ed2a4857f3
% canonical-lessons: ES-C715-amar, ES-C715-detestar, ES-C715-enamorado, ES-C715-besar, ES-C715-pasion

\chapter{Amar y detestar}
\label{ch:es-amar-y-detestar}
\hlchaptermodality{715}

\begin{chapteropening}
\textbf{By the end of this chapter:} \emph{I can talk about love, passion and loathing in a story or a song.}

Complete the last lesson of chapter 715: I can talk about love, passion and loathing in a story or a song.
\end{chapteropening}

\section[amar]{amar — to love}

\label{lesson:ES-C715-amar}



\begin{quote}
Before the new one: say the Spanish for tender, then the Spanish for to adore.
\end{quote}

\subsection*{You'll want to know: amar}
\textbf{amar} — \textquotedblleft{}to love\textquotedblright{}.

\begin{quote}
Te amo con todo mi corazón.
\end{quote}

\begin{quote}
\emph{I love you with all my heart.}
\end{quote}

Deep, romantic love; with friends and family people usually say \emph{te quiero}.

\begin{cousinweb}
\emph{amar} is inherited from Latin \emph{amāre} \textquotedblleft{}to love\textquotedblright{} (EtymDB); English \emph{amateur}, one who does something for love, comes from it through French.
\end{cousinweb}

\begin{grammarlens}[title={What you've built}]
One more word for telling, explaining and describing at B1.
\end{grammarlens}

\subsection*{Your turn}
\begin{itemize}
\raggedright
  \item \emph{Say it:} \emph{amar}
  \item \emph{Say it:} \emph{amar}, once more
  \item \emph{Say it:} a sentence of your own with \emph{amar}
  \item \emph{Recall:} say the Spanish for tender, then the Spanish for to adore, then say \emph{amar} again
\end{itemize}

\begin{tcolorbox}[breakable,colback=teal!4,colframe=teal!35!black,title={Before you move on}]
What does amar mean? (\textquotedblleft{}To love\textquotedblright{}.) Say it once more.
\end{tcolorbox}
\section[detestar]{detestar — to detest, to loathe}

\label{lesson:ES-C715-detestar}



\begin{quote}
Before the new one: say the Spanish for to adore, then the Spanish for to love.
\end{quote}

\subsection*{You'll want to know: detestar}
\textbf{detestar} — \textquotedblleft{}to detest, to loathe\textquotedblright{}.

\begin{quote}
Detesto las mentiras.
\end{quote}

\begin{quote}
\emph{I loathe lies.}
\end{quote}

Close to \emph{odiar}, and a little more formal. It takes a direct object: \emph{detesto el frío}.

\begin{cousinweb}
\emph{detestar} is a learned borrowing from Latin \emph{dētestārī} \textquotedblleft{}to curse, to call the gods to witness against\textquotedblright{}, built on \emph{testis} \textquotedblleft{}witness\textquotedblright{}.
\end{cousinweb}

\begin{grammarlens}[title={What you've built}]
One more word for telling, explaining and describing at B1.
\end{grammarlens}

\subsection*{Your turn}
\begin{itemize}
\raggedright
  \item \emph{Say it:} \emph{detestar}
  \item \emph{Say it:} \emph{detestar}, once more
  \item \emph{Say it:} a sentence of your own with \emph{detestar}
  \item \emph{Recall:} say the Spanish for to adore, then the Spanish for to love, then say \emph{detestar} again
\end{itemize}

\begin{tcolorbox}[breakable,colback=teal!4,colframe=teal!35!black,title={Before you move on}]
What does detestar mean? (\textquotedblleft{}To detest, to loathe\textquotedblright{}.) Say it once more.
\end{tcolorbox}
\section[enamorado]{enamorado — in love}

\label{lesson:ES-C715-enamorado}



\begin{quote}
Before the new one: say the Spanish for to love, then the Spanish for to detest.
\end{quote}

\subsection*{You'll want to know: enamorado}
\textbf{enamorado} — \textquotedblleft{}in love\textquotedblright{}.

\begin{quote}
Están muy enamorados.
\end{quote}

\begin{quote}
\emph{They are very much in love.}
\end{quote}

\emph{Estar enamorado de alguien} — to be in love with someone.

\begin{cousinweb}
\emph{enamorado} is built on \emph{enamorar} \textquotedblleft{}to make fall in love\textquotedblright{}, from \emph{en-} and \emph{amor} \textquotedblleft{}love\textquotedblright{} (EtymDB); \emph{amor} comes from Latin \emph{amor}.
\end{cousinweb}

\begin{grammarlens}[title={What you've built}]
One more word for telling, explaining and describing at B1.
\end{grammarlens}

\subsection*{Your turn}
\begin{itemize}
\raggedright
  \item \emph{Say it:} \emph{enamorado}
  \item \emph{Say it:} \emph{enamorado}, once more
  \item \emph{Say it:} a sentence of your own with \emph{enamorado}
  \item \emph{Recall:} say the Spanish for to love, then the Spanish for to detest, then say \emph{enamorado} again
\end{itemize}

\begin{tcolorbox}[breakable,colback=teal!4,colframe=teal!35!black,title={Before you move on}]
What does enamorado mean? (\textquotedblleft{}In love\textquotedblright{}.) Say it once more.
\end{tcolorbox}
\section[besar]{besar — to kiss}

\label{lesson:ES-C715-besar}



\begin{quote}
Before the new one: say the Spanish for to detest, then the Spanish for in love.
\end{quote}

\subsection*{You'll want to know: besar}
\textbf{besar} — \textquotedblleft{}to kiss\textquotedblright{}.

\begin{quote}
Me besó en la cara.
\end{quote}

\begin{quote}
\emph{She kissed me on the face.}
\end{quote}

\emph{Un beso} is a kiss; \emph{besos} closes many messages.

\begin{cousinweb}
\emph{besar} is usually traced to Latin \emph{bāsiāre} \textquotedblleft{}to kiss\textquotedblright{} (Lewis \& Short), built on \emph{bāsium} \textquotedblleft{}a kiss\textquotedblright{}; the Latin \emph{-si-} became \emph{-s-}.
\end{cousinweb}

\begin{grammarlens}[title={What you've built}]
One more word for telling, explaining and describing at B1.
\end{grammarlens}

\subsection*{Your turn}
\begin{itemize}
\raggedright
  \item \emph{Say it:} \emph{besar}
  \item \emph{Say it:} \emph{besar}, once more
  \item \emph{Say it:} a sentence of your own with \emph{besar}
  \item \emph{Recall:} say the Spanish for to detest, then the Spanish for in love, then say \emph{besar} again
\end{itemize}

\begin{tcolorbox}[breakable,colback=teal!4,colframe=teal!35!black,title={Before you move on}]
What does besar mean? (\textquotedblleft{}To kiss\textquotedblright{}.) Say it once more.
\end{tcolorbox}
\section[la pasión]{la pasión — passion}

\label{lesson:ES-C715-pasion}



\begin{quote}
Before the new one: say the Spanish for in love, then the Spanish for to kiss.
\end{quote}

\subsection*{You'll want to know: la pasión}
\textbf{la pasión} — \textquotedblleft{}passion\textquotedblright{}.

\begin{quote}
Baila con pasión.
\end{quote}

\begin{quote}
\emph{She dances with passion.}
\end{quote}

\emph{Tener pasión por algo} — to be passionate about something.

\begin{cousinweb}
\emph{pasión} is a learned borrowing from Latin \emph{passiō} \textquotedblleft{}suffering\textquotedblright{} (EtymDB), built on \emph{patī} \textquotedblleft{}to suffer\textquotedblright{}; strong feeling was something you suffered.
\end{cousinweb}

\begin{grammarlens}[title={What you've built}]
One more word for telling, explaining and describing at B1.
\end{grammarlens}

\subsection*{Your turn}
\begin{itemize}
\raggedright
  \item \emph{Say it:} \emph{la pasión}
  \item \emph{Say it:} \emph{la pasión}, once more
  \item \emph{Say it:} a sentence of your own with \emph{la pasión}
  \item \emph{Recall:} say the Spanish for in love, then the Spanish for to kiss, then say \emph{la pasión} again
\end{itemize}

\begin{tcolorbox}[breakable,colback=teal!4,colframe=teal!35!black,title={Before you move on}]
What does la pasión mean? (\textquotedblleft{}Passion\textquotedblright{}.) Say it once more.
\end{tcolorbox}
